\chapter{Plan de pruebas}
\label{chap:plan_pruebas}

En este capítulo se abordará todo el proceso relacionado con el plan de pruebas, incluyendo las pruebas realizadas durante el desarrollo, las pruebas posteriores, quiénes han participado en ellas y cómo se han identificado y solucionado los errores detectados. Además, se detallarán los problemas surgidos a lo largo del desarrollo de la plataforma, cómo se resolvieron y si afectaron al progreso del proyecto llegando a provocar retrasos en otras tareas.

En primer lugar, es importante destacar que las pruebas se llevaron a cabo de forma continua durante el proceso de desarrollo e implementación. No se avanzaba a nuevas tareas sin haber resuelto previamente todos los errores detectados en la etapa actual. Por ejemplo, si se observaba que el sistema de autenticación no propagaba correctamente el token JWT a las peticiones subsiguientes, no se procedía con la siguiente funcionalidad, como podría ser la vista del dashboard, hasta que el problema estuviera resuelto. No obstante, en algunos casos los errores se detectaron una vez iniciada una nueva tarea, debido a que solo se manifestaban bajo condiciones muy específicas que no se habían dado antes. En estos casos, se debía decidir entre terminar la tarea actual si estaba próxima a completarse o interrumpirla para solucionar el error anterior, lo que inevitablemente provocaba ciertos retrasos en el desarrollo.

En segundo lugar, una vez completada toda la aplicación, se llevaron a cabo numerosas pruebas para verificar que todas las funcionalidades y datos se comportaban como se deseaba.

En tercer lugar, tanto el sistema como la aplicación fueron probados por diferentes personas: amigos, compañeros y, en especial, el tutor, con el fin de recoger opiniones, detectar posibles mejoras en la experiencia de usuario y observar cómo interactuaban con la plataforma. Esto permitió identificar errores adicionales que fueron corregidos, así como realizar mejoras visuales que ayudaron a clarificar ciertas funcionalidades que no estaban siendo bien percibidas, y también aportaron ideas de posibles implementaciones que en algunos casos se valoraron y añadieron al proyecto, mientras que otras se dejaron como posibles mejoras futuras.

Una vez tratados estos aspectos generales, el capítulo se dividirá en tres secciones: en la primera se expondrán los principales problemas surgidos durante el desarrollo; en la segunda se explicará cómo se abordó y resolvió cada uno de ellos; y en la tercera se expondrán los diferentes entornos en los que se realizaron las pruebas de la plataforma, tanto a nivel de dispositivos como de navegadores.

\section{Problemas durante el desarrollo}
\label{sec:problemas_desarrollo}

En esta sección se comentan los problemas, dificultades y fallos que fueron surgiendo a lo largo del desarrollo de Colivi.

La \textbf{Prueba-01} consistió en establecer la comunicación entre el servidor de desarrollo del frontend, ejecutándose en el puerto 5173 gracias a Vite, y la API REST del backend en el puerto 8080. Desde el primer momento, el navegador bloqueaba todas las peticiones salientes del cliente React con errores de política de mismo origen (\emph{Same-Origin Policy}), impidiendo cualquier interacción con el servidor hasta que se configurara correctamente CORS.

\begin{table}[H]
\centering
\begin{tabular}{|p{0.25\textwidth}|p{0.65\textwidth}|}
\hline
\textbf{Prueba-01} & \textbf{CORS entre frontend y backend} \\ \hline
Salida esperada & Comunicación fluida entre el cliente React (puerto 5173) y la API REST (puerto 8080). \\ \hline
Salida obtenida & Error en consola: Access to XMLHttpRequest blocked by CORS policy. Ninguna petición HTTP llegaba al backend. \\ \hline
\end{tabular}
\caption{Tabla Prueba-01}
\label{tab:prueba01}
\end{table}

En cuanto a la \textbf{Prueba-02}, una vez que se iniciaron las migraciones de base de datos con Flyway, se detectó que cualquier modificación sobre un script de migración ya ejecutado provocaba que la aplicación no arrancara. Flyway almacena un \emph{checksum} de cada archivo aplicado y, al detectar que el contenido había cambiado respecto al valor registrado, lanzaba una excepción bloqueante que impedía el inicio del servidor.

\begin{table}[H]
\centering
\begin{tabular}{|p{0.25\textwidth}|p{0.65\textwidth}|}
\hline
\textbf{Prueba-02} & \textbf{Integridad de migraciones Flyway} \\ \hline
Salida esperada & La aplicación arranca con normalidad tras corregir un script de migración. \\ \hline
Salida obtenida & La aplicación no arranca: FlywayException: Validate failed — Migration checksum mismatch. \\ \hline
\end{tabular}
\caption{Tabla Prueba-02}
\label{tab:prueba02}
\end{table}

Otro de los problemas importantes fue el que se refiere a la \textbf{Prueba-03}, relacionado con las consultas a la base de datos al cargar el panel principal (\emph{dashboard}). Al inspeccionar los logs de Hibernate en modo DEBUG, se observaba que por cada miembro del hogar se lanzaba una consulta adicional para obtener sus datos relacionados, multiplicando innecesariamente la carga sobre PostgreSQL y aumentando los tiempos de respuesta de forma proporcional al número de convivientes.

\begin{table}[H]
\centering
\begin{tabular}{|p{0.25\textwidth}|p{0.65\textwidth}|}
\hline
\textbf{Prueba-03} & \textbf{Problema N+1 en Hibernate} \\ \hline
Salida esperada & El resumen del hogar contiene todos los datos necesarios sin requerir consultas adicionales. \\ \hline
Salida obtenida & Hibernate ejecuta 1 + N consultas adicionales (una por cada miembro y relación lazy), degradando el rendimiento. \\ \hline
\end{tabular}
\caption{Tabla Prueba-03}
\label{tab:prueba03}
\end{table}

También se detectaron dificultades con el envío de correos transaccionales, como se puede observar en la \textbf{Prueba-04}. Los flujos de recuperación de contraseña y reactivación de cuenta requerían el envío de correos con enlaces temporales, pero el servidor de Spring Boot rechazaba la autenticación SMTP contra Gmail, impidiendo el despacho de cualquier mensaje.

\begin{table}[H]
\centering
\begin{tabular}{|p{0.25\textwidth}|p{0.65\textwidth}|}
\hline
\textbf{Prueba-04} & \textbf{Autenticación SMTP para correos transaccionales} \\ \hline
Salida esperada & El usuario recibe un correo con el enlace de recuperación o reactivación de cuenta. \\ \hline
Salida obtenida & Error en el servidor: 535-5.7.8 Username and Password not accepted. No se envía ningún correo. \\ \hline
\end{tabular}
\caption{Tabla Prueba-04}
\label{tab:prueba04}
\end{table}

La visualización y actualización reactiva de datos en el frontend (\textbf{Prueba-05}) fue uno de los problemas principales. Al realizar operaciones de escritura —como registrar un nuevo gasto o completar una tarea doméstica—, la interfaz no reflejaba los cambios de forma inmediata. El estado global gestionado con Redux Toolkit no se sincronizaba con los datos recién obtenidos del servidor, obligando al usuario a recargar la página para ver el resultado actualizado.

\begin{table}[H]
\centering
\begin{tabular}{|p{0.25\textwidth}|p{0.65\textwidth}|}
\hline
\textbf{Prueba-05} & \textbf{Sincronización del estado tras mutaciones} \\ \hline
Entrada & Registrar un nuevo gasto compartido o marcar una tarea como completada. \\ \hline
Salida esperada & La vista se actualiza automáticamente reflejando el cambio sin recargar la página. \\ \hline
Salida obtenida & La vista permanece desactualizada hasta que el usuario recarga manualmente la página. \\ \hline
\end{tabular}
\caption{Tabla Prueba-05}
\label{tab:prueba05}
\end{table}

Otro obstáculo fue la integración del mapa interactivo con Leaflet en la vista de búsqueda de alojamientos (\textbf{Prueba-06}). Leaflet accede directamente al objeto window del navegador, que no existe durante el proceso de compilación de Vite, generando errores que impedían construir el proyecto.

\begin{table}[H]
\centering
\begin{tabular}{|p{0.25\textwidth}|p{0.65\textwidth}|}
\hline
\textbf{Prueba-06} & \textbf{Compatibilidad de Leaflet con el entorno de compilación} \\ \hline
Salida esperada & El mapa interactivo se renderiza correctamente en el explorador de alojamientos. \\ \hline
Salida obtenida & Error de compilación: ReferenceError: window is not defined. El proyecto no podía construirse. \\ \hline
\end{tabular}
\caption{Tabla Prueba-06}
\label{tab:prueba06}
\end{table}

La \textbf{Prueba-07} consistió en la implementación de la geocodificación de direcciones mediante la API pública de Nominatim. Al publicar un alojamiento, la aplicación convertía la dirección postal a coordenadas geográficas; sin embargo, el servicio externo impone un límite estricto de una petición por segundo y bloquea las IPs que superan dicha cuota, lo que generaba errores intermitentes durante las pruebas con múltiples alojamientos.

\begin{table}[H]
\centering
\begin{tabular}{|p{0.25\textwidth}|p{0.65\textwidth}|}
\hline
\textbf{Prueba-07} & \textbf{Rate limiting en la API de geocodificación Nominatim} \\ \hline
Entrada & Publicación simultánea o rápida de varios alojamientos con nuevas direcciones. \\ \hline
Salida esperada & Todas las direcciones son geocodificadas y los alojamientos se muestran en el mapa. \\ \hline
Salida obtenida & Error HTTP 429 Too Many Requests. Las coordenadas no se almacenan y los alojamientos no aparecen en el mapa. \\ \hline
\end{tabular}
\caption{Tabla Prueba-07}
\label{tab:prueba07}
\end{table}

También surgieron complicaciones con la definición de las herramientas del asistente de inteligencia artificial (\textbf{Prueba-08}). Spring AI expone las herramientas MCP mediante esquemas JSON derivados de las clases Java, pero cualquier discrepancia entre los tipos del esquema y los DTOs provocaba respuestas malformadas o errores de deserialización silenciosos que eran muy difíciles de reproducir y depurar.

\begin{table}[H]
\centering
\begin{tabular}{|p{0.25\textwidth}|p{0.65\textwidth}|}
\hline
\textbf{Prueba-08} & \textbf{Contrato de tipos en herramientas MCP de Spring AI} \\ \hline
Salida esperada & El asistente de IA invoca correctamente las herramientas y devuelve respuestas coherentes. \\ \hline
Salida obtenida & El modelo devuelve respuestas malformadas o vacías; errores de deserialización silenciosos en el servidor. \\ \hline
\end{tabular}
\caption{Tabla Prueba-08}
\label{tab:prueba08}
\end{table}

El primero de los problemas relacionados con la versión móvil generada con Capacitor fue que nunca se había trabajado previamente con esta tecnología en un proyecto de esta envergadura (\textbf{Prueba-09}). En dispositivos físicos Android, las peticiones al backend en localhost:8080 fallaban sistemáticamente, ya que el dispositivo no comparte la interfaz de red loopback del equipo de desarrollo.

\begin{table}[H]
\centering
\begin{tabular}{|p{0.25\textwidth}|p{0.65\textwidth}|}
\hline
\textbf{Prueba-09} & \textbf{Conectividad del backend en dispositivos físicos Android} \\ \hline
Salida esperada & La aplicación móvil se comunica con el backend y muestra los datos correctamente. \\ \hline
Salida obtenida & Todas las peticiones fallan con error de red; el dispositivo no puede resolver localhost:8080. \\ \hline
\end{tabular}
\caption{Tabla Prueba-09}
\label{tab:prueba09}
\end{table}


\section{Solución de los problemas encontrados}
\label{sec:soluciones_problemas}

En este apartado se expone cómo se han ido resolviendo los distintos problemas que surgieron en el desarrollo, descritos en el apartado anterior.

Por un lado, respecto a la \textbf{Prueba-01} relacionada con el bloqueo de peticiones por la política de mismo origen, se implementó una configuración global de CORS mediante un bean de tipo WebMvcConfigurer anotado con @Configuration. Este bean define un CorsConfigurationSource que autoriza explícitamente el origen http://localhost:5173, los métodos HTTP requeridos (GET, POST, PUT, DELETE, PATCH, OPTIONS) y las cabeceras necesarias (Authorization, Content-Type), aplicándose de forma global a todos los endpoints bajo /api/v1/**.

\begin{table}[H]
\centering
\begin{tabular}{|p{0.25\textwidth}|p{0.65\textwidth}|}
\hline
\textbf{Prueba-01} & \textbf{CORS entre frontend y backend} \\ \hline
Salida esperada & Comunicación fluida entre el cliente React (puerto 5173) y la API REST (puerto 8080). \\ \hline
Salida obtenida & Error en consola: Access to XMLHttpRequest blocked by CORS policy. \\ \hline
Solución & Implementar un bean WebMvcConfigurer con CorsConfigurationSource que autorice explícitamente el origen del frontend y los métodos HTTP requeridos. \\ \hline
\end{tabular}
\caption{Tabla Solución Prueba-01}
\label{tab:sol_prueba01}
\end{table}

En relación con la \textbf{Prueba-02} sobre la integridad de las migraciones Flyway, la solución adoptada fue respetar rigurosamente la inmutabilidad de los scripts ya aplicados. Ante cualquier cambio en el esquema, se creaba siempre un nuevo script con la versión incremental siguiente (por ejemplo, V3__add_column_x.sql). Para los entornos de desarrollo donde era necesario reiniciar el esquema desde cero, se utilizó la propiedad flyway.cleanDisabled=false, que permite limpiar y reaplicar todas las migraciones de forma controlada sin afectar a otros entornos.

\begin{table}[H]
\centering
\begin{tabular}{|p{0.25\textwidth}|p{0.65\textwidth}|}
\hline
\textbf{Prueba-02} & \textbf{Integridad de migraciones Flyway} \\ \hline
Salida esperada & La aplicación arranca con normalidad tras corregir un script de migración. \\ \hline
Salida obtenida & La aplicación no arranca: FlywayException: Validate failed --- Migration checksum mismatch. \\ \hline
Solución & No modificar nunca scripts ya aplicados; crear siempre un nuevo script de versión incremental para cualquier cambio en el esquema. \\ \hline
\end{tabular}
\caption{Tabla Solución Prueba-02}
\label{tab:sol_prueba02}
\end{table}

Respecto al problema N+1 de Hibernate de la \textbf{Prueba-03}, la solución consistió en diseñar consultas JPQL personalizadas en los repositorios con cláusulas LEFT JOIN FETCH sobre las relaciones involucradas. Para el endpoint consolidado del dashboard, se implementó un DTO de proyección (HomeSummaryDTO) que permite al servidor agregar todos los datos necesarios ---miembros, gastos y tareas--- en un número acotado de consultas SQL, eliminando por completo las consultas adicionales de Hibernate.

\begin{table}[H]
\centering
\begin{tabular}{|p{0.25\textwidth}|p{0.65\textwidth}|}
\hline
\textbf{Prueba-03} & \textbf{Problema N+1 en Hibernate} \\ \hline
Salida esperada & El endpoint resuelve todos los datos del hogar en una sola consulta SQL. \\ \hline
Salida obtenida & Hibernate ejecuta 1 + N consultas adicionales, degradando el rendimiento. \\ \hline
Solución & Consultas JPQL con LEFT JOIN FETCH y DTO de proyección HomeSummaryDTO para agregar todos los datos en una única transacción de lectura. \\ \hline
\end{tabular}
\caption{Tabla Solución Prueba-03}
\label{tab:sol_prueba03}
\end{table}

En cuanto a la \textbf{Prueba-04} del servicio de correo, la solución definitiva fue sustituir la autenticación estándar de Gmail por una \textbf{contraseña de aplicación} generada desde la configuración de seguridad de la cuenta de Google, combinada con la habilitación de la verificación en dos pasos. Las credenciales se externalizaron al archivo application.properties mediante variables de entorno, siguiendo el principio de no incluir secretos en el repositorio de código.

\begin{table}[H]
\centering
\begin{tabular}{|p{0.25\textwidth}|p{0.65\textwidth}|}
\hline
\textbf{Prueba-04} & \textbf{Autenticación SMTP para correos transaccionales} \\ \hline
Salida esperada & El usuario recibe un correo con el enlace de recuperación o reactivación de cuenta. \\ \hline
Salida obtenida & Error en el servidor: 535-5.7.8 Username and Password not accepted. \\ \hline
Solución & Generar una contraseña de aplicación específica en la cuenta de Google y configurarla en application.properties como variable de entorno. \\ \hline
\end{tabular}
\caption{Tabla Solución Prueba-04}
\label{tab:sol_prueba04}
\end{table}

La solución adoptada para la \textbf{Prueba-05}, relativa a la falta de sincronización del estado tras mutaciones, fue migrar la gestión del estado asíncrono de Redux Toolkit a \textbf{TanStack Query}. Las peticiones de lectura se encapsulan en useQuery y las mutaciones en useMutation, con invalidación reactiva de claves de caché (queryClient.invalidateQueries) tras cada operación de escritura. Esta arquitectura redujo significativamente el código de gestión de estado y garantizó la coherencia de los datos entre componentes sin intervención manual del usuario.

\begin{table}[H]
\centering
\begin{tabular}{|p{0.25\textwidth}|p{0.65\textwidth}|}
\hline
\textbf{Prueba-05} & \textbf{Sincronización del estado tras mutaciones} \\ \hline
Entrada & Registrar un nuevo gasto compartido o marcar una tarea como completada. \\ \hline
Salida esperada & La vista se actualiza automáticamente reflejando el cambio. \\ \hline
Salida obtenida & La vista permanece desactualizada hasta recargar la página. \\ \hline
Solución & Sustituir Redux Toolkit por TanStack Query con invalidación reactiva de claves de caché tras cada mutación. \\ \hline
\end{tabular}
\caption{Tabla Solución Prueba-05}
\label{tab:sol_prueba05}
\end{table}

Para resolver el error de compilación de la \textbf{Prueba-06} causado por Leaflet, el componente del mapa se aisló mediante carga dinámica (\emph{lazy loading}) con React.lazy y Suspense. De este modo, Leaflet solo se inicializa en el contexto del navegador y nunca durante el proceso de compilación de Vite, eliminando por completo el acceso prematuro al objeto window.

\begin{table}[H]
\centering
\begin{tabular}{|p{0.25\textwidth}|p{0.65\textwidth}|}
\hline
\textbf{Prueba-06} & \textbf{Compatibilidad de Leaflet con el entorno de compilación} \\ \hline
Salida esperada & El mapa interactivo se renderiza correctamente. \\ \hline
Salida obtenida & Error de compilación: ReferenceError: window is not defined. \\ \hline
Solución & Envolver el componente del mapa en React.lazy + Suspense para garantizar que Leaflet solo se ejecute en el navegador. \\ \hline
\end{tabular}
\caption{Tabla Solución Prueba-06}
\label{tab:sol_prueba06}
\end{table}

Respecto al \emph{rate limiting} de Nominatim de la \textbf{Prueba-07}, la solución consistió en implementar una estrategia de \textbf{geocodificación asíncrona con persistencia local}: la llamada a Nominatim se realiza una única vez en el momento en que el anfitrión publica o actualiza un alojamiento, y las coordenadas resultantes se almacenan permanentemente en las columnas latitude y longitude de la tabla accommodations en PostgreSQL. Las búsquedas habituales del cliente consultan directamente la base de datos local sin realizar ninguna llamada externa.

\begin{table}[H]
\centering
\begin{tabular}{|p{0.25\textwidth}|p{0.65\textwidth}|}
\hline
\textbf{Prueba-07} & \textbf{Rate limiting en la API de geocodificación Nominatim} \\ \hline
Entrada & Publicación simultánea de varios alojamientos con nuevas direcciones. \\ \hline
Salida esperada & Todas las direcciones son geocodificadas y los alojamientos aparecen en el mapa. \\ \hline
Salida obtenida & Error HTTP 429 Too Many Requests; coordenadas no almacenadas. \\ \hline
Solución & Geocodificar una única vez al publicar el alojamiento y persistir las coordenadas en base de datos, eliminando las llamadas externas en las búsquedas. \\ \hline
\end{tabular}
\caption{Tabla Solución Prueba-07}
\label{tab:sol_prueba07}
\end{table}

Para resolver las incompatibilidades de la \textbf{Prueba-08} con Spring AI, se adoptó una estrategia de tipado estricto bidireccional: los esquemas JSON de las herramientas se generaron directamente desde las clases Java mediante la anotación @Tool de Spring AI, garantizando que el contrato de datos fuese siempre coherente con la implementación. Se habilitó el modo de depuración de Spring AI en los logs para visualizar las invocaciones del modelo y detectar desajustes de forma temprana.

\begin{table}[H]
\centering
\begin{tabular}{|p{0.25\textwidth}|p{0.65\textwidth}|}
\hline
\textbf{Prueba-08} & \textbf{Contrato de tipos en herramientas MCP de Spring AI} \\ \hline
Salida esperada & El asistente de IA invoca correctamente las herramientas y devuelve respuestas coherentes. \\ \hline
Salida obtenida & Respuestas malformadas o vacías; errores de deserialización silenciosos. \\ \hline
Solución & Generar los esquemas JSON directamente desde las clases Java con @Tool y activar el modo DEBUG de Spring AI para auditar las invocaciones. \\ \hline
\end{tabular}
\caption{Tabla Solución Prueba-08}
\label{tab:sol_prueba08}
\end{table}

Por último, en relación con la \textbf{Prueba-09} de conectividad en dispositivos físicos Android, la solución fue configurar la URL base del backend mediante la IP local del equipo de desarrollo (por ejemplo, http://192.168.1.X:8080) en el archivo .env específico para el entorno móvil. Adicionalmente, se habilitó una regla de entrada en el firewall de Windows para permitir conexiones al puerto 8080 desde dispositivos en la misma red local.

\begin{table}[H]
\centering
\begin{tabular}{|p{0.25\textwidth}|p{0.65\textwidth}|}
\hline
\textbf{Prueba-09} & \textbf{Conectividad del backend en dispositivos físicos Android} \\ \hline
Salida esperada & La aplicación móvil se comunica con el backend y muestra los datos correctamente. \\ \hline
Salida obtenida & Todas las peticiones fallan con error de red; el dispositivo no puede resolver localhost:8080. \\ \hline
Solución & Sustituir localhost por la IP local del equipo en el archivo .env móvil y abrir el puerto 8080 en el firewall de Windows para la red local. \\ \hline
\end{tabular}
\caption{Tabla Solución Prueba-09}
\label{tab:sol_prueba09}
\end{table}


\section{Entornos de pruebas y dispositivos utilizados}
\label{sec:entornos_pruebas}

En este apartado se describen los entornos y dispositivos empleados para verificar el correcto funcionamiento de la plataforma, tanto en su versión web como en su adaptación para dispositivos móviles.

\subsection{Entorno de desarrollo}

El desarrollo de Colivi se llevó a cabo en un único equipo de escritorio. Las herramientas principales utilizadas fueron las siguientes:

\begin{itemize}
    \item \textbf{Sistema operativo}: Windows 11 Pro.
    \item \textbf{IDE principal}: IntelliJ IDEA Ultimate para el desarrollo del backend en Spring Boot y Java.
    \item \textbf{Editor de código}: Visual Studio Code para el desarrollo del frontend en React y TypeScript.
    \item \textbf{Gestión de base de datos}: DBeaver, utilizado para inspeccionar el esquema de PostgreSQL, ejecutar consultas SQL de diagnóstico y verificar el estado de las migraciones de Flyway.
    \item \textbf{Pruebas de API}: Postman, empleado para verificar los endpoints REST del backend de forma independiente al frontend, comprobando códigos de respuesta HTTP y contratos de los DTOs.
    \item \textbf{Control de versiones}: Git con repositorio remoto en GitHub.
\end{itemize}

\subsection{Navegadores web utilizados}

La versión web de Colivi fue probada en los siguientes navegadores de escritorio:

\begin{itemize}
    \item \textbf{Google Chrome}: navegador principal durante el desarrollo. Se emplearon intensivamente las herramientas para desarrolladores (DevTools), especialmente el panel \emph{Network} para auditar las peticiones HTTP, el panel \emph{Application} para inspeccionar el localStorage y los tokens JWT, y la consola JavaScript para depurar el comportamiento de Leaflet, Supercluster y TanStack Query.
    \item \textbf{Mozilla Firefox}: utilizado para verificar la compatibilidad entre navegadores, especialmente en lo relativo al comportamiento del mapa interactivo y las animaciones CSS.
    \item \textbf{Microsoft Edge}: empleado para pruebas adicionales de compatibilidad en entorno Windows.
\end{itemize}

\subsection{Dispositivos móviles}

Para la validación de la versión móvil generada con Capacitor, se realizaron pruebas en dispositivos físicos con sistema operativo Android, conectados mediante depuración USB (\emph{ADB}). Esto permitió inspeccionar los logs de la aplicación en tiempo real desde las DevTools de Chrome mediante la opción chrome://inspect, verificando el renderizado de los componentes, la navegación entre vistas, el funcionamiento del mapa y la comunicación con el backend a través de la red local.

\subsection{Integración continua con GitHub Actions}

Como complemento a las pruebas manuales, se configuró un flujo de integración continua (\emph{CI}) mediante \textbf{GitHub Actions}. Cada vez que se realizaba un \emph{push} o se abría una \emph{pull request} sobre la rama principal del repositorio, se disparaba automáticamente un \emph{workflow} definido en el archivo .github/workflows/ci.yml que ejecutaba las siguientes comprobaciones:

\begin{itemize}
    \item \textbf{Compilación del backend}: el flujo levantaba un contenedor con Java 21 y Maven, compilaba el proyecto de Spring Boot y ejecutaba los tests unitarios e de integración con mvn verify. Cualquier fallo en la compilación o en los tests bloqueaba la integración del cambio.
    \item \textbf{Compilación del frontend}: en paralelo, otro \emph{job} instalaba las dependencias de Node.js y ejecutaba npm run build sobre el proyecto Vite + React, verificando que el código TypeScript compilaba sin errores de tipo y que el empaquetado producía los artefactos estáticos correctamente.
    \item \textbf{Notificación de resultados}: el estado de cada ejecución quedaba visible directamente en GitHub junto a cada \emph{commit}, permitiendo identificar de forma inmediata qué cambio introdujo una regresión sin necesidad de ejecutar las comprobaciones de forma manual en el equipo local.
\end{itemize}

Este sistema de CI actuó como una red de seguridad automatizada a lo largo de todo el desarrollo, garantizando que la rama principal del repositorio se mantuviese siempre en un estado compilable y funcional.

\subsection{Despliegue en la nube con Microsoft Azure}

Para poder realizar pruebas en un entorno real de producción, así como para que la aplicación fuese accesible desde cualquier dispositivo sin depender del equipo de desarrollo local, se aprovisionó una \textbf{máquina virtual en Microsoft Azure}. Esta infraestructura en la nube permitió verificar el comportamiento de la plataforma bajo condiciones reales de red, alejadas del entorno localhost.

La máquina virtual se configuró con los siguientes componentes:

\begin{itemize}
    \item \textbf{Sistema operativo}: Ubuntu Server LTS, elegido por su estabilidad y por la amplia disponibilidad de documentación para el despliegue de aplicaciones Java y Node.js.
    \item \textbf{Backend}: el servidor de Spring Boot se desplegó como un servicio del sistema (\emph{systemd}), garantizando que el proceso se reiniciase automáticamente tras un reinicio de la máquina o un fallo inesperado.
    \item \textbf{Base de datos}: instancia de PostgreSQL instalada en la misma máquina virtual, con acceso restringido únicamente al proceso del backend mediante la configuración del firewall de Azure (\emph{Network Security Group}).
    \item \textbf{Frontend}: los artefactos estáticos generados por Vite (npm run build) se sirvieron mediante \textbf{Nginx}, configurado como proxy inverso para redirigir las peticiones /api/** al backend en el puerto 8080 y servir el cliente React en el puerto 80/443.
    \item \textbf{Certificado SSL}: se configuró un certificado TLS mediante \textbf{Certbot} y \emph{Let's Encrypt}, habilitando HTTPS en la aplicación, requisito indispensable para que la versión móvil generada con Capacitor pudiera comunicarse con el servidor de forma segura.
\end{itemize}

Este entorno de Azure sirvió también como plataforma de pruebas de aceptación final, permitiendo que personas externas al desarrollo (amigos, compañeros y el tutor) accediesen a la aplicación desde sus propios dispositivos y navegadores sin ninguna configuración adicional.

